\documentclass[a4paper, serif, 10pt]{moderncv}

%% moderncv
% color and style
\moderncvcolor{black}
\moderncvstyle{banking}

%% page layout/margins
\usepackage[top=2.5cm, bottom=2.5cm, left=1.5cm, right=1.5cm]{geometry}

\nopagenumbers{}

%% Babel config for Brazilian Portuguese language
% ref:
% - https://latex3.github.io/babel/guides/locale-portuguese.html
\usepackage[brazilian]{babel}

%% fontspec
\usepackage{fontspec}

\IfFontExistsTF{Linux Libertine}{
  \setmainfont[Ligatures={TeX, Common}, Numbers=OldStyle]{Linux Libertine}
}{}
\IfFontExistsTF{Linux Libertine O}{
  \setmainfont[Ligatures={TeX, Common}, Numbers=OldStyle]{Linux Libertine O}
}{}

%% CTeX
% CJK support, used to show CJK characters in the resume
%
% - heading=false: disable ctex heading
% - punct=kaiming: use kaiming punctuations styles for CJK
% - scheme=plain: use plain scheme, do not override \normalsize font size
% - space=auto: space settings for CJK characters
%
% CJK fonts are resolved per locale language, e.g. Noto Serif/Sans CJK JP
% for Japanese resumes, with fallback to the next region if not installed
%
% ref:
% - http://ctan.mirrorcatalogs.com/language/chinese/ctex/ctex.pdf
\usepackage[UTF8, heading=false, punct=kaiming, scheme=plain, space=auto]{ctex}

\IfFontExistsTF{Noto Serif CJK SC}{
  \setCJKmainfont{Noto Serif CJK SC}
}{}
\IfFontExistsTF{Noto Sans CJK SC}{
  \setCJKsansfont{Noto Sans CJK SC}
}{}

%% line spacing
\usepackage{setspace}
\setstretch{1.125}

%% URL styling - use normal text instead of monospace
\urlstyle{same}

%% Patch \httplink and \httpslink to handle full URLs
% The original moderncv macros always prepend http:// or https:// to URLs.
% This causes issues when the URL already has a protocol (e.g., https://example.com)
% resulting in malformed URLs like https://https://example.com.
% This patch checks if the URL already contains "://" and uses it directly if so.
% Note: We use \str_set:Nx (x-expansion) to expand macros like \@homepage first.
\ExplSyntaxOn
\str_new:N \g_yamlresume_url_str

\RenewDocumentCommand{\httpslink}{O{}m}{
  \str_set:Nx \g_yamlresume_url_str {#2}
  \str_if_in:NnTF \g_yamlresume_url_str {://}
    { \url{#2} }
    { \url{https://#2} }
}

\RenewDocumentCommand{\httplink}{O{}m}{
  \str_set:Nx \g_yamlresume_url_str {#2}
  \str_if_in:NnTF \g_yamlresume_url_str {://}
    { \url{#2} }
    { \url{http://#2} }
}
\ExplSyntaxOff

\name{Mariana Costa Almeida}{}
\title{Gerente de Operações | Logística, S\&OP e Melhoria Contínua}
\phone[mobile]{+55 11 98765-4321}
\email{mariana.almeida@example.com}
\homepage{https://www.marianaalmeida.com.br}

\address{Rua das Acácias, 245, São Paulo, São Paulo, Brasil, 04567-000}{}{}

\extrainfo{{\small \faLinkedin}\ \href{https://www.linkedin.com/in/mariana-almeida-ops}{@mariana-almeida-ops} {} {} {} • {} {} {} 
{\small \faMedium}\ \href{https://medium.com/@marianaalmeida-ops}{@marianaalmeida-ops}}

\begin{document}

\maketitle

\section{Informações Básicas}

\cvline{}{Gerente de Operações com mais de 10 anos de experiência em logística, planejamento e melhoria contínua. Histórico comprovado na liderança de equipes multidisciplinares, otimização de processos e redução de custos operacionais em empresas de médio e grande porte.
%
\begin{itemize}
\item Especialista em S\&OP, gestão de estoques e indicadores de desempenho (KPIs)
\item Experiência com Lean Six Sigma, mapeamento de processos e gestão de mudanças
\item Forte atuação em negociação com fornecedores, transportadoras e parceiros estratégicos
\item Perfil analítico, orientado a dados e focado em resultados sustentáveis
\end{itemize}}
\section{Formação}

\cventry{fev. de 2010–dez. de 2014}
        {Bacharelado, Engenharia de Produção, Nota: 8.7/10}
        {Universidade de São Paulo}
        {\url{https://www.usp.br}}
        {}
        {\begin{itemize}
\item Formação com foco em engenharia de operações, logística e melhoria de processos
\item Projeto de conclusão sobre otimização de rotas de distribuição
\item Participação em empresa júnior e projetos de consultoria operacional
\end{itemize}
\textbf{Cursos}: Gestão de Operações, Logística e Cadeia de Suprimentos, Planejamento e Controle da Produção, Pesquisa Operacional, Gestão da Qualidade, Engenharia Econômica, Simulação de Sistemas, Gestão de Projetos}

\cventry{jan. de 2016–dez. de 2017}
        {Mestrado, Gestão de Operações e Cadeia de Suprimentos, Nota: 9.2/10}
        {Fundação Getulio Vargas}
        {\url{https://www.fgv.br}}
        {}
        {\begin{itemize}
\item Especialização em gestão estratégica de operações e cadeia de suprimentos
\item Desenvolvimento de projeto aplicado à redução de custos logísticos
\item Ênfase em liderança, negociação e transformação digital
\end{itemize}}
\section{Experiência Profissional}

\cventry{mar. de 2021–Atual}
        {Gerente de Operações}
        {LogiRota Soluções Logísticas}
        {\url{https://www.logirota.com.br}}
        {}
        {\begin{itemize}
\item Lidero equipe de 45 pessoas nas áreas de transporte, armazenagem e planejamento operacional
\item Reduzi o custo logístico em 18\% por meio de renegociação de contratos e otimização de rotas
\item Implementei rotina de S\&OP integrada às áreas comercial, compras e produção
\item Estruturei painéis de KPIs para acompanhamento diário de nível de serviço, custo por entrega e produtividade
\item Conduzi projeto de automação de processos com integração entre WMS e TMS
\end{itemize}
\textbf{Palavras-chave}: Gestão de operações, Logística, S\&OP, Lean, KPIs, Liderança}

\cventry{jan. de 2018–fev. de 2021}
        {Coordenadora de Operações}
        {Distribuidora Nova Era}
        {\url{https://www.distribuidoranovaera.com.br}}
        {}
        {\begin{itemize}
\item Coordenei operações de distribuição em três centros de distribuição regionais
\item Reduzi o lead time médio de entrega em 22\% com revisão de processos e roteirização
\item Implantei programa 5S e rotinas de auditoria operacional nas unidades
\item Gerenciei orçamento anual de R\$ 12 milhões, com acompanhamento mensal de desvios
\item Desenvolvi indicadores de acuracidade de estoque e nível de atendimento ao cliente
\end{itemize}
\textbf{Palavras-chave}: Distribuição, Estoques, 5S, Melhoria contínua, Gestão de equipes}

\cventry{jul. de 2014–dez. de 2017}
        {Analista de Operações Sênior}
        {Transportes Bandeirante}
        {\url{https://www.transportesbandeirante.com.br}}
        {}
        {\begin{itemize}
\item Analisei indicadores operacionais e elaborei relatórios gerenciais para a diretoria
\item Desenvolvi previsões de demanda e apoiei o planejamento de capacidade da frota
\item Apoiei a implantação do sistema WMS e treinamento das equipes operacionais
\item Participei de projetos de redução de custos com transporte e consolidação de cargas
\item Realizei negociações com fornecedores e transportadoras parceiras
\end{itemize}
\textbf{Palavras-chave}: Análise de dados, Previsão de demanda, WMS, Transporte}
\section{Idiomas}

\cvline{Português}{Nativo ou Bilíngue}
\cvline{Inglês}{Proficiência Profissional Fluente \hfill \textbf{Palavras-chave}: TOEFL 108, Reuniões internacionais}
\cvline{Espanhol}{Proficiência Profissional Limitada \hfill \textbf{Palavras-chave}: DELE B2}
\section{Competências}

\cvline{Gestão de Operações}{Especialista \hfill \textbf{Palavras-chave}: Planejamento de demanda, S\&OP, Lean Six Sigma, Gestão de indicadores, Redução de custos}
\cvline{Logística e Distribuição}{Avançado \hfill \textbf{Palavras-chave}: Roteirização, WMS, TMS, Armazenagem, Gestão de estoques}
\cvline{Análise de Dados}{Avançado \hfill \textbf{Palavras-chave}: Excel avançado, Power BI, SQL, Dashboards, Análise de indicadores}
\cvline{Liderança e Pessoas}{Especialista \hfill \textbf{Palavras-chave}: Gestão de equipes, Coaching, Cultura de segurança, Desenvolvimento de talentos}
\cvline{Gestão de Projetos}{Avançado \hfill \textbf{Palavras-chave}: PMBOK, Scrum, Melhoria contínua, Gestão de mudanças}
\section{Prêmios}

\cventry{nov. de 2023}
        {LogiRota Soluções Logísticas}
        {Prêmio de Excelência Operacional}
        {}
        {}
        {Reconhecimento pela liderança do projeto de centralização do controle logístico, com ganho de 15\% em produtividade e melhoria consistente no nível de serviço ao cliente.}

\cventry{dez. de 2020}
        {Distribuidora Nova Era}
        {Reconhecimento por Redução de Custos}
        {}
        {}
        {Premiação pelo resultado de redução de 22\% no lead time de entrega e melhoria da acuracidade de estoque nas operações regionais.}
\section{Certificados}

\cventry{jun. de 2019}
        {ASQ}
        {Lean Six Sigma Green Belt}
        {\url{https://asq.org/cert/lean-six-sigma-green-belt}}
        {}
        {}

\cventry{set. de 2020}
        {ASCM}
        {Certified Supply Chain Professional (CSCP)}
        {\url{https://www.ascm.org/credentials-certifications/ascm-certifications/cscp/}}
        {}
        {}

\cventry{mar. de 2022}
        {Scrum Alliance}
        {Scrum Master Certified}
        {\url{https://www.scrumalliance.org/get-certified/scrum-master-track/certified-scrummaster}}
        {}
        {}
\section{Publicações}

\cventry{ago. de 2022}
        {Análise preditiva aplicada à gestão de estoques no varejo}
        {Revista Brasileira de Logística}
        {\url{https://www.revistabrasileiradelogistica.com.br/artigos/analise-preditiva-estoques}}
        {}
        {\begin{itemize}
\item Estudo sobre o uso de modelos preditivos para melhorar a acuracidade de estoque
\item Discussão de impactos na redução de rupturas e no capital de giro
\item Apresentação de casos práticos em operações de distribuição
\end{itemize}}
\section{Projetos}

\cventry{jan. de 2022–jul. de 2022}
        {Projeto de centralização do monitoramento de entregas e indicadores operacionais.}
        {Centro de Controle Logístico}
        {}
        {}
        {\begin{itemize}
\item Implantação de sala de controle com acompanhamento em tempo real das operações
\item Integração de dados de WMS, TMS e atendimento ao cliente
\item Redução de 12\% no tempo de resposta a ocorrências logísticas
\end{itemize}
\textbf{Palavras-chave}: Logística, Controle operacional, KPIs, Integração de sistemas}

\cventry{mar. de 2019–nov. de 2019}
        {Programa de melhoria contínua aplicado aos processos de separação e expedição.}
        {Programa Lean na Distribuição}
        {}
        {}
        {\begin{itemize}
\item Mapeamento de fluxo de valor e identificação de desperdícios
\item Implantação de padrões operacionais e rotinas de gestão à vista
\item Ganho de 20\% em produtividade nas equipes de separação
\end{itemize}
\textbf{Palavras-chave}: Lean, Melhoria contínua, Produtividade, Gestão à vista}
\section{Interesses}

\cvline{Esportes}{Corrida, Ciclismo, Vôlei}
\cvline{Leitura}{Gestão, Biografias, Economia}
\cvline{Culinária}{Cozinha italiana, Panificação}
\section{Voluntariado}

\cventry{fev. de 2020–Atual}
        {Mentora de Jovens em Logística}
        {Instituto Rede Solidária}
        {\url{https://www.institutoredesolidaria.org.br}}
        {}
        {\begin{itemize}
\item Mentoria para jovens em início de carreira nas áreas de logística e operações
\item Apoio na elaboração de currículos e preparação para entrevistas
\item Participação em workshops sobre planejamento de carreira e competências técnicas
\item Colaboração com a organização de feiras de empregabilidade comunitárias
\end{itemize}}
    
\end{document}